% wave17_a2_preface_rectification_plan.tex
% Wave 17, A2/20. Positive-geometry grammar for Y^+(X) in
% chapters/frame/preface.tex. Plan only; no preface edits.
\providecommand{\Phid}{\Phi_d}
\providecommand{\PhiFA}{\Phi^{\mathrm{FA}}}
\providecommand{\SpCh}{\mathrm{Sp}^{\mathrm{ch}}}
\providecommand{\CoHA}{\mathrm{CoHA}}
\providecommand{\Heq}{H^{*}_{\mathrm{eq}}}
\providecommand{\Meff}{\mathcal{M}^{+}_{\mathrm{eff}}}
\providecommand{\CY}{\mathrm{CY}}
\providecommand{\kch}{\kappa_{\mathrm{ch}}}

\section*{Cycle 1 --- Foundational $\Phi$ statement (lines 1--250)}

The boxed $\Phi_d$ (lines 28--38) and the (U1)--(U4) block
(lines 54--88) state the two-stage factorisation at the level
of holomorphic factorisation algebras on~$X$ and their
specialisations. The (U4) \emph{standard-input recovery}
discharges the positive half case-by-case: $d{=}1$ lattice
VOA; $d{=}2$ Mukai Heisenberg of rank $24$; $d{=}3$ route
$\PhiFA_3(\CoHA(\C^3)) \leadsto Y^{+}(\widehat{\frakgl}_1)$
via an $\Omega$-backgrounded $\C^2$-specialisation (lines
83--87). The wording treats $\CoHA(\C^3)$ as categorical input
and names Schiffmann--Vasserot at the end of the clause; the
cascade table (lines 950--984) repeats this by phenotype
(K3 fibre, Nakajima cycle, pt).

\textbf{Defect.} No universal grammar for the positive half
is asserted. The toric $\Einf$/VOA case, the K3 Mukai case,
and the generic CoHA case each read as local recoveries; the
reader sees four faces (lattice VOA / Mukai Heisenberg /
affine Yangian / $\CoHA(Q_X,W_X)$) without the single rule
that subsumes them.

\textbf{Proposed target.} After line 88, before the `four-step
recipe' paragraph, insert a short anchor sentence of the form
\[
 Y^{+}(X)
 \;=\;
 \Heq\!\bigl(\Meff(X)\bigr),
\]
with the equivariance type declared by geometric class
(toric $\Rightarrow$ full $T$; compact reduced CY$_3$
$\Rightarrow$ reduced equivariance after DT/reduced-obstruction;
orbifold $\Rightarrow$ discrete isotropy; polarised
$\Rightarrow$ period-domain $\mathrm{Sp}/\mathrm{O}$-action).
Kontsevich--Soibelman 2008 (critical CoHA, full-$T$
equivariance on the moduli of stable objects), Van den Bergh
2004 (quiver-with-potential as the categorical datum),
Schiffmann--Vasserot 2013 (the isomorphism $\CoHA(Q_{\C^3}) \cong
Y^{+}(\widehat{\frakgl}_1)$ via the torus-equivariant push on
$\mathrm{Hilb}^n(\C^2)$), Oberdieck--Pixton 2017 (reduced DT on
$K3 \times E$ and the reduced equivariance of its moduli).
Two sentences; no narration.

\section*{Cycle 2 --- The $G(X)$ clause (lines 500--750)}

Lines 579--597 (Mukai lattice on the Lie-theoretic side) state
the positive-half identification
\[
 \CoHA(Q_X, W_X)
 \;\cong\;
 Y^{+}\!\bigl(\widehat{\frakg}_X\bigr),
\]
citing Schiffmann--Vasserot and
Rap\v{c}\'ak--Soibelman--Yang--Zhao, then unpacks the Jordan
case and the K3 Mukai case. Lines 715--720 state
$C(\frakg, q) = D(Y^{+}(\frakg_{K3}))$ at the abelian level
and flag the non-abelian full Drinfeld double as Conjecture
CY-C. Lines 788--870 (pentagon block) enumerate six routes
with generator ranks $\rho^{R_i} \in \{3, 12, 24\}$.

\textbf{Defect.} The passage from ``Jordan quiver'' to ``K3
Mukai quiver'' to ``$K3 \times E$ pentagon'' happens by
enumeration. The reduced-DT equivariance specific to
compact reduced CY$_3$ (Oberdieck--Pixton, removing the
translation-of-$E$ zero mode) is not visible in the preface.
The non-toric generator count $\rho \in \{3, 12, 24\}$ is
tabulated but not pegged to an equivariance type.

\textbf{Proposed target.} Before the pentagon block (line
787), insert a two-sentence bridge: the equivariance class of
$\Meff(X)$ stratifies $Y^{+}(X)$ and pins each pentagon node
to its equivariance datum --- $R_1$ uses reduced DT on
$K3 \times E$ (Oberdieck--Pixton $\mathrm{II}_{1,1}$-descent),
$R_3$ uses the full Narain torus on $\Lambda^{3,2}$, $R_4$
uses the $\Z/2$-Kummer isotropy, $R_5,R_6$ use the K3
Mukai-$\mathrm{O}(4,20)$ period-domain equivariance.
The generator rank $\rho^{R_i}$ is then the rank of the
equivariance-fixed Cartan; the pentagon is a diagram of
equivariance-type reductions, not six unrelated quantum-group
constructions.

\section*{Cycle 3 --- The universal sentence}

Insert after line 88 (foundational), mirrored at line 787
(pentagon), one paragraph:

\begin{quote}
Across $\CY_d$-$\mathrm{Cat}$, the positive half is the
equivariant cohomology of the positive effective moduli,
\[
 Y^{+}(X)
 \;=\;
 \Heq\!\bigl(\Meff(X)\bigr),
\]
where the equivariance type stratifies by geometric class:
full torus~$T$ on toric CY$_d$ (Kontsevich--Soibelman 2008);
reduced equivariance after the $\mathrm{II}_{1,1}$-descent on
compact reduced CY$_3$ carrying a non-trivial BPS zero mode
(Oberdieck--Pixton 2017); discrete isotropy on orbifold
quotients (Van den Bergh 2004 quiver-with-potential picture);
period-domain $\mathrm{O}(2, n{+}2)$-equivariance on polarised
CY$_d$ through the Mukai lattice action. The Schiffmann--Vasserot
identification $\CoHA(Q_X, W_X) \cong Y^{+}(\widehat{\frakg}_X)$
is the rank-unstable evaluation of the universal formula on
the toric fan; the reduced-equivariance evaluation at
$K3 \times E$ is the source of the five pentagon nodes; the
Drinfeld double $G(X) = D(Y^{+}(X))$ is the uniform closure,
with $G(K3) = Y(\frakg_{K3})$ the ADE/Kummer-chart face
and $G(K3 \times E)$ the pentagon colimit when the
equivariance types of $R_1, R_3, R_4, R_5, R_6$ agree after
reduced-equivariant transport.
\end{quote}

One paragraph. Equals signs, not analogies. Equivariance
type, not enumeration. Nine lines of prose, zero bookkeeping
vocabulary, CG voice.

\section*{Cycle 4 --- Weaving into Drinfeld-double and two-stage threads}

\textbf{Drinfeld-double thread} (lines 66--73, 170--171,
1874--1877, 1924--1927): each passage reads
$\cZ(\Rep^{\Eone}(A)) \simeq \Rep^{\Etwo}(Z^{\mathrm{der}}(A))$
and identifies $Y^{+}$ as the ordered-Hall input. The
Cycle-3 paragraph makes $Y^{+}$ the same object across all
four passages, with equivariance-type as the additional datum.
No prose edit is required in lines 66--73 or 170--171 if the
Cycle-3 paragraph precedes them: the identification
$Y^{+}(X) = \Heq(\Meff(X))$ is then referenced by symbol only.

\textbf{Two-stage thread} (lines 926--985, cascade table):
the $\PhiFA_d$ column becomes the pair (positive effective
moduli, equivariance type); the $\SpCh$ column is unchanged;
the shadow column inherits the universal $Y^{+}$-formula via
the bar complex (ordered Hall $\Rightarrow$ bar differential
of $B^{\mathrm{ord}}(A_\cC)$; Vol I companion paper). One
footnote on the table's $\PhiFA_d$ column cell for $d = 3$,
adding ``(reduced equivariance for compact non-toric; full
torus for toric CY$_3$)'', is sufficient to weave Cycle 3
into the cascade. No other edit to lines 950--985.

\textbf{Pentagon thread} (lines 787--870): the equivariance
reading (Cycle 2) replaces the pentagon's ``six different
constructions'' framing by ``five equivariance-type reductions
of one positive-effective-moduli cohomology plus one Borcherds
source''. $\rho^{R_i}$ becomes the equivariance-fixed Cartan
rank.

\section*{Cycle 5 --- Plan, not patch}

No prose is to be inscribed by this wave. Targets:
(i) insert after line 88 (Cycle 3 paragraph);
(ii) insert before line 787 (Cycle 2 bridge, two sentences);
(iii) annotate the cascade table cell at $d{=}3$ (Cycle 4
footnote). Citations: Kontsevich--Soibelman 2008
(arXiv:0811.2435); Van den Bergh 2004 (\emph{A relation
between Hochschild homology and cohomology for Gorenstein
rings}, Proc.\ AMS, reformulated here as the
quiver-with-potential datum $(Q_X, W_X)$ in the
non-commutative-crepant-resolution line 1865 of the preface);
Schiffmann--Vasserot 2013 (\emph{Cherednik algebras,
$\mathcal{W}$-algebras and the equivariant cohomology of the
moduli space of instantons on $\mathbf{A}^2$},
arXiv:1202.2756); Oberdieck--Pixton 2017 (\emph{Holomorphic
anomaly equations and the Igusa cusp form conjecture},
arXiv:1706.10100, reduced-DT zero-mode on $K3 \times E$).
No fabrication; every citation is already in the preface's
ambient citation field. Execution waits on authorisation;
this plan is the atom.
